\chapter{مفاهیم پایه و مرور کارهای پیشین}
\label{ch:background}

این فصل دو کار انجام می‌دهد. نخست، مفاهیم و ابزارهایی را معرفی می‌کند که فهم فصل‌های بعد به
آن‌ها وابسته است: \lr{Kubernetes} و الگوریتم دقیق مقیاس‌گذار پیش‌فرض آن، \lr{Prometheus} و
زبان پرس‌وجوی آن، روش‌های پیش‌بینی سری زمانی، و مفاهیم لازم از کنترل پس‌خوردی. دوم، کارهای
پیشین را مرور می‌کند و نشان می‌دهد شکافی که این پروژه در آن قرار می‌گیرد کجاست.

معرفی مفاهیم در این فصل گزینشی است و نه دانشنامه‌ای. تنها آن بخشی از هر ابزار شرح داده می‌شود
که در فصل‌های بعد به آن ارجاع داده خواهد شد. برای نمونه، از \lr{Kubernetes} تنها منابع مربوط
به مقیاس‌گذاری و چرخه‌ی آماده‌شدن یک نمونه توضیح داده می‌شود، چون همان‌ها هستند که در
اندازه‌گیری‌های فصل~\ref{ch:results} نقش دارند.

\section{رایانش ابری و معماری ریزخدمت}
\label{sec:bg-cloud}

رایانش ابری منابع پردازشی را به‌صورت خدمتی برحسب تقاضا در اختیار می‌گذارد. ویژگی‌ای از آن که
برای این نوشتار اهمیت دارد، امکان تغییر سریع مقدار منابع در اختیار است.

معماری ریزخدمت سامانه را به سرویس‌های کوچک، مستقل و جداگانه‌مستقرشدنی می‌شکند. سه پیامد این
معماری برای مقیاس‌گذاری خودکار اهمیت دارد:

\begin{itemize}
  \item \textbf{واحد مقیاس‌گذاری کوچک‌تر می‌شود.} به‌جای افزودن یک نسخه از کل سامانه، تنها
  نمونه‌های سرویسی افزوده می‌شود که زیر فشار است.
  \item \textbf{سرویس‌های بدون‌حالت}\LTRfootnote{Stateless} رایج‌اند، و افزودن یا حذف
  نمونه‌ی آن‌ها نیازمند جابه‌جایی داده نیست. این پروژه تنها چنین سرویس‌هایی را در نظر می‌گیرد.
  \item \textbf{سنجه‌های سطح‌برنامه در دسترس‌اند.} هر سرویس نرخ درخواست، تأخیر و خطای خود را
  منتشر می‌کند، که سیگنال‌های مستقیم‌تری از بار نسبت به مصرف پردازنده‌اند.
\end{itemize}

\section{Kubernetes}
\label{sec:bg-k8s}

\lr{Kubernetes} یک سکوی ارکستراسیون کانتینر است. چند منبع آن در این نوشتار به‌طور مکرر ظاهر
می‌شوند و در ادامه معرفی می‌شوند.

\subsection{منابع پایه}
\label{sec:bg-k8s-resources}

\textbf{\lr{Pod}} کوچک‌ترین واحد استقرار است و یک یا چند کانتینر را در بر می‌گیرد. در این
نوشتار، هر \lr{Pod} دقیقاً یک نمونه از سرویس هدف است و واژه‌ی «نمونه» به همین معنا به کار
می‌رود.

\textbf{\lr{Deployment}} حالت مطلوب مجموعه‌ای از \lr{Pod}‌های یکسان را توصیف می‌کند. میدان
\lr{spec.replicas} می‌گوید چند نمونه \emph{خواسته شده} است و میدان
\lr{status.readyReplicas} می‌گوید چند نمونه در حال حاضر \emph{آماده‌ی سرویس} است. تفاوت این
دو عدد، که در حالت پایدار صفر است، در طول یک مقیاس‌گیری به بالا مثبت می‌شود. همان‌طور که در
بخش~\ref{sec:design-formula} خواهیم دید، انتخاب اینکه کدام‌یک پایه‌ی فرمول باشد، یک تصمیم
طراحی با پیامد عددی مشخص است.

\textbf{\lr{Service}} یک نقطه‌ی دسترسی پایدار برای مجموعه‌ای از \lr{Pod}‌ها فراهم می‌کند و
ترافیک را میان نمونه‌های آماده توزیع می‌کند. توزیع ترافیک در سطح \emph{اتصال} انجام می‌شود و
نه در سطح \emph{درخواست}؛ این جزئیات ظاهراً کم‌اهمیت، در بخش~\ref{sec:results-loaddefect} به
یک نقص جدی اندازه‌گیری می‌انجامد.

\textbf{زیرمنبع \lr{scale}} یک واسط استاندارد برای خواندن و نوشتن تعداد نمونه‌های یک منبع
مقیاس‌پذیر است. هر کنترل‌گر مقیاس‌گذاری، از جمله \lr{HPA} و مقیاس‌گذار این پروژه، از همین
واسط استفاده می‌کند.

\subsection{چرخه‌ی آماده‌شدن یک نمونه}
\label{sec:bg-k8s-readiness}

این بخش برای فهم زمان واکنش در فصل~\ref{ch:results} ضروری است. از لحظه‌ای که تعداد نمونه‌ها
افزایش می‌یابد تا لحظه‌ای که نمونه‌ی تازه واقعاً ترافیک می‌گیرد، این مراحل طی می‌شود:
زمان‌بندی روی یک گره، دریافت تصویر کانتینر (اگر محلی نباشد)، آغاز فرایند، و سرانجام موفق شدن
کاوشگر آمادگی\LTRfootnote{Readiness Probe}. تنها پس از موفقیت کاوشگر آمادگی است که
\lr{Service} ترافیک را به آن نمونه می‌فرستد و \lr{status.readyReplicas} افزایش می‌یابد.

دو کاوشگر در این نوشتار نقش دارند:

\begin{itemize}
  \item \textbf{کاوشگر آمادگی} تعیین می‌کند نمونه ترافیک بگیرد یا نه. چون همین شمارش، پایه‌ی
  فرمول مقیاس‌گذار است، تنظیم این کاوشگر مستقیماً بر تصمیم‌های مقیاس‌گذاری اثر می‌گذارد.
  \item \textbf{کاوشگر زنده‌بودن}\LTRfootnote{Liveness Probe} تعیین می‌کند نمونه بازراه‌اندازی
  شود یا نه. اصل مهمی که در بخش~\ref{sec:method-observability-collapse} با شواهد تجربی به آن
  بازمی‌گردیم این است که کاوشگر زنده‌بودن باید یک فرایند \emph{قفل‌شده} را تشخیص دهد و نه یک
  فرایند \emph{کُند}؛ بازراه‌اندازی یک فرایند پرمشغله تقریباً هیچ‌گاه پاسخ درستی نیست.
\end{itemize}

در بستر آزمایش این پروژه، فاصله‌ی میان صدور فرمان مقیاس‌گیری تا آماده‌شدن نمونه حدود ۳۰ ثانیه
است. این عدد در کنار فاصله‌ی آشتی‌دهی ۱۵ ثانیه‌ای، مقیاس زمانی‌ای را می‌سازد که کل بحث
پیش‌بینی در این نوشتار حول آن می‌چرخد: افق پیش‌بینی باید تقریباً هم‌اندازه‌ی همین تأخیر باشد
تا ظرفیت به‌موقع آماده شود.

\section{مقیاس‌گذاری خودکار در Kubernetes}
\label{sec:bg-autoscaling-k8s}

\subsection{مقیاس‌گذار افقی پیش‌فرض و الگوریتم آن}
\label{sec:bg-hpa}

\lr{Horizontal Pod Autoscaler} به‌صورت دوره‌ای — به‌طور پیش‌فرض هر ۱۵ ثانیه — سنجه‌ها را
می‌خواند و تعداد نمونه‌های مطلوب را از رابطه‌ی زیر حساب می‌کند~\cite{k8sdocs-hpa}:

\begin{equation}
\label{eq:hpa}
r_{\text{desired}} = \left\lceil r_{\text{current}} \times
\frac{m_{\text{current}}}{m_{\text{target}}} \right\rceil
\end{equation}

که در آن $m_{\text{current}}$ مقدار جاری سنجه به ازای هر نمونه و $m_{\text{target}}$ مقدار
هدف است. سه ویژگی این الگوریتم در سراسر این نوشتار به کار می‌آید.

\textbf{ناحیه‌ی مرده.} اگر نسبت $m_{\text{current}}/m_{\text{target}}$ به اندازه‌ی کافی به یک
نزدیک باشد، هیچ تغییری اعمال نمی‌شود. آستانه‌ی پیش‌فرض ۱۰ درصد است. هدف این ناحیه، جلوگیری از
تعقیب نوسان‌های کوچک سنجه است.

\textbf{پنجره‌ی پایدارسازی و سیاست‌های \lr{behavior}.} در نسخه‌ی \lr{autoscaling/v2}،
میدان \lr{behavior} امکان تنظیم جداگانه‌ی رفتار مقیاس به بالا و مقیاس به پایین را می‌دهد.
مهم‌ترین تنظیم آن \lr{stabilizationWindowSeconds} است: به‌جای عمل کردن بر پایه‌ی آخرین
پیشنهاد، کنترل‌گر پیشنهادها را در طول یک پنجره نگه می‌دارد و از آن‌ها بیشینه (برای مقیاس به
پایین) می‌گیرد. به این ترتیب یک افت کوتاه بار نمی‌تواند ناوگان را کوچک کند. مقدار پیش‌فرض این
پنجره برای مقیاس به پایین به‌طور محسوسی بزرگ‌تر از مقیاس به بالاست، که بازتاب یک نامتقارنی
واقعی است: هزینه‌ی دیر افزودن ظرفیت، نقض کیفیت خدمت است، اما هزینه‌ی دیر کاستن ظرفیت، فقط
پول است.

\textbf{قاعده‌ی سنجه‌های غایب.}\LTRfootnote{Missing Metrics} اگر سنجه‌ی برخی نمونه‌ها در
دسترس نباشد، \lr{HPA} محافظه‌کارانه عمل می‌کند: در محاسبه‌ی مقیاس به \emph{پایین} فرض می‌کند
آن نمونه‌ها در ۱۰۰ درصد مقدار هدف‌اند، و در محاسبه‌ی مقیاس به \emph{بالا} فرض می‌کند در صفر
درصدند. نمونه‌های هنوز آماده‌نشده نیز از میانگین‌گیری کنار گذاشته می‌شوند. این قاعده در
بخش~\ref{sec:results-hpacustom-redo} نقش تعیین‌کننده‌ای پیدا می‌کند: در حضور یک نقص در تحویل
بار، همین قاعده مقیاس به پایین بازوی \lr{hpa-custom} را به‌طور کامل مسدود کرده بود.

\subsection{مقیاس‌گذاری عمودی و مقیاس‌گذاری خوشه}
\label{sec:bg-vpa-ca}

دو سازوکار دیگر نیز وجود دارند که در این پروژه به کار نرفته‌اند اما برای کامل بودن تصویر
معرفی می‌شوند. \lr{Vertical Pod Autoscaler} به‌جای تعداد نمونه‌ها، درخواست منابع هر نمونه را
تنظیم می‌کند. \lr{Cluster Autoscaler} تعداد گره‌های خوشه را تغییر می‌دهد تا نمونه‌های
زمان‌بندی‌نشده جا پیدا کنند. هر سه سازوکار در عمل با هم به کار می‌روند و تعامل میان آن‌ها خود
موضوع پژوهش مستقلی است که خارج از محدوده‌ی این پروژه است.

\subsection{سنجه‌های سفارشی: از metrics-server تا Prometheus Adapter}
\label{sec:bg-custom-metrics}

\lr{HPA} در ساده‌ترین پیکربندی خود، سنجه‌های منبعی\LTRfootnote{Resource Metrics} را از
\lr{metrics-server} می‌خواند؛ یعنی مصرف پردازنده و حافظه. برای سنجه‌های سطح‌برنامه،
\lr{Kubernetes} سه واسط تعریف کرده است: \lr{Custom Metrics API}،
\lr{External Metrics API} و \lr{Resource Metrics API}. این واسط‌ها خودشان پیاده‌سازی ندارند و
باید توسط یک مؤلفه‌ی بیرونی ارائه شوند.

\lr{Prometheus Adapter} یکی از پیاده‌سازی‌های رایج این واسط است~\cite{prometheusadapter}. کار
آن، ترجمه‌ی پرس‌وجوهای \lr{PromQL} به پاسخ‌های \lr{Custom Metrics API} است. در این پروژه از
همین ابزار برای ساختن بازوی مقایسه‌ی \lr{hpa-custom} استفاده شده است، که در آن \lr{HPA} روی
\emph{همان سنجه‌ای} کار می‌کند که سیاست ما استفاده می‌کند. اهمیت این بازو در
بخش~\ref{sec:method-arms} توضیح داده می‌شود: بدون آن، مقایسه میان دو تغییر هم‌زمان (سیگنال
بهتر و سیاست بهتر) مخدوش می‌ماند.

نکته‌ای که در فصل~\ref{ch:results} به آن بازمی‌گردیم: سنجه‌های نوع \lr{Pods} در این واسط
\emph{به ازای هر نمونه} گزارش می‌شوند. این یعنی بازوی \lr{hpa-custom} تنها بازویی است که
ورودی‌اش به تفکیک نمونه است، و در نتیجه تنها بازویی است که نسبت به توزیع نایکنواخت ترافیک
میان نمونه‌ها آسیب‌پذیر است.

\subsection{KEDA}
\label{sec:bg-keda}

\lr{KEDA} یک مقیاس‌گذار رویدادمحور\LTRfootnote{Event-driven} برای \lr{Kubernetes} است که در
سال ۲۰۲۳ به وضعیت پروژه‌ی فارغ‌التحصیل بنیاد \lr{CNCF} رسید~\cite{kedagraduation}. این ابزار
مقیاس‌گذاری بر پایه‌ی منابع بیرونی — از جمله \lr{Prometheus}، صف‌ها و
\lr{Kafka} — و همچنین مقیاس‌گذاری تا صفر\LTRfootnote{Scale-to-zero} را فراهم می‌کند. در عمل،
\lr{KEDA} یک \lr{HPA} را مدیریت می‌کند و خودش سیاست تصمیم‌گیری تازه‌ای نمی‌آورد.

وجود \lr{KEDA} دقیقاً همان دلیلی است که در بخش~\ref{sec:intro-question} گفته شد: ساختن یک
مقیاس‌گذار مبتنی بر \lr{Prometheus} کار حل‌شده‌ای است. آنچه \lr{KEDA} ارائه نمی‌کند، تحلیل
پایداری سیاست پیش‌بینانه است.

\section{Prometheus و زبان پرس‌وجوی آن}
\label{sec:bg-prometheus}

\lr{Prometheus} یک سامانه‌ی پایش با مدل کشیدن\LTRfootnote{Pull} است: به‌صورت دوره‌ای نقاط
انتهایی \lr{/metrics} را می‌خواند و مقادیر را به‌عنوان سری زمانی ذخیره می‌کند~\cite{prometheus}.
هر سری با یک نام سنجه و مجموعه‌ای از برچسب‌ها مشخص می‌شود.

سه نوع سنجه در این نوشتار به کار رفته است:

\begin{itemize}
  \item \textbf{شمارنده}\LTRfootnote{Counter}: مقداری که تنها افزایش می‌یابد، مانند
  \lr{http\_requests\_total}. مقدار خام یک شمارنده معنای مستقیمی ندارد؛ آنچه معنا دارد نرخ
  تغییر آن است.
  \item \textbf{سنجه‌ی لحظه‌ای}\LTRfootnote{Gauge}: مقداری که بالا و پایین می‌رود، مانند
  تعداد درخواست‌های در جریان یا سنجه‌های داخلی خود مقیاس‌گذار.
  \item \textbf{بافه}\LTRfootnote{Histogram}: توزیع مقادیر در سطل‌های از پیش تعیین‌شده، که
  برای محاسبه‌ی صدک‌های تأخیر به کار می‌رود.
\end{itemize}

سیگنال مقیاس‌گذاری در این پروژه با عبارت زیر خوانده می‌شود:

\begin{latin}
\begin{lstlisting}[language=promql]
sum(rate(http_requests_total{app="sample"}[1m]))
  / clamp_min(count(up{app="sample"} == 1), 1)
\end{lstlisting}
\end{latin}

\subsection{تأخیر ذاتی مسیر اندازه‌گیری}
\label{sec:bg-measurement-lag}

نکته‌ای که برای فهم نتایج این پروژه اهمیت زیادی دارد، تأخیر انباشته در مسیر اندازه‌گیری است.
از لحظه‌ای که بار واقعاً تغییر می‌کند تا لحظه‌ای که مقیاس‌گذار آن را می‌بیند، سه تأخیر جمع
می‌شود:

\begin{enumerate}
  \item \textbf{فاصله‌ی جمع‌آوری.} \lr{Prometheus} هر چند ثانیه یک بار نمونه‌برداری می‌کند
  (در این پروژه ۵ ثانیه برای سنجه‌های مهم).
  \item \textbf{پنجره‌ی \lr{rate}.} تابع \lr{rate(...[1m])} شیب را روی یک پنجره‌ی یک‌دقیقه‌ای
  حساب می‌کند، پس در عمل یک فیلتر پایین‌گذر با تأخیر گروهی حدود نیمِ پنجره است.
  \item \textbf{فاصله‌ی آشتی‌دهی.} مقیاس‌گذار هر ۱۵ ثانیه یک بار تصمیم می‌گیرد.
\end{enumerate}

مجموع این تأخیرها ده‌ها ثانیه است، یعنی هم‌مرتبه با زمان آماده‌شدن یک نمونه. این نکته دو
پیامد دارد که در فصل~\ref{ch:results} به هر دو بازمی‌گردیم: نخست آنکه بخش قابل توجهی از سودی
که پیش‌بینی می‌آورد، در واقع \emph{جبران همین تأخیر اندازه‌گیری} است و نه پیش‌بینی آینده به
معنای دقیق کلمه؛ و دوم آنکه شبیه‌سازی که این تأخیر را مدل نکند، به‌طور نظام‌مند سیاست واکنشی
را بهتر از آنچه هست نشان می‌دهد (بخش~\ref{sec:results-simulator}).

\section{پیش‌بینی سری زمانی}
\label{sec:bg-forecasting}

هدف پیش‌بینی‌کننده در این پروژه ساده است: با در دست داشتن دنباله‌ای از مقادیر گذشته‌ی بار،
مقدار آن را برای $h$ گام بعد تخمین بزند.

\subsection{هموارسازی نمایی}
\label{sec:bg-ewma}

ساده‌ترین روش کارآمد، میانگین متحرک با وزن نمایی\LTRfootnote{Exponentially Weighted Moving
Average (EWMA)} است:

\begin{equation}
\label{eq:ewma}
\ell_k = \alpha\, x_k + (1-\alpha)\, \ell_{k-1}
\end{equation}

که در آن $x_k$ مشاهده‌ی گام $k$ و $\alpha \in (0,1]$ ضریب هموارسازی است. مقادیر بزرگ‌تر
$\alpha$ وزن بیشتری به مشاهده‌های تازه می‌دهند و پاسخ را سریع‌تر اما پرنوسان‌تر می‌کنند.

\lr{EWMA} به‌تنهایی \emph{روند} را دنبال نمی‌کند: اگر سری به‌طور پیوسته بالا برود، خروجی
هموارشده همیشه عقب می‌ماند. برای برون‌یابی به آینده باید شیب هم برآورد شود. ساده‌ترین شکل آن،
که در این پروژه با نام \lr{EWMATrend} پیاده‌سازی شده، برآورد خطی شیب از مقادیر هموارشده و
سپس برون‌یابی $h$ گام به جلو است.

\subsection{روش Holt}
\label{sec:bg-holt}

هولت~\cite{holt2004} شکل کامل‌تری از همین ایده را پیشنهاد کرد که سطح و روند را با دو ضریب
جداگانه هموار می‌کند:

\begin{align}
\label{eq:holt}
\ell_k &= \alpha\, x_k + (1-\alpha)\,(\ell_{k-1} + b_{k-1}) \\
b_k &= \beta\,(\ell_k - \ell_{k-1}) + (1-\beta)\, b_{k-1} \\
\hat{x}_{k+h} &= \ell_k + h\, b_k
\end{align}

وینترز~\cite{winters1960} بعدها مؤلفه‌ی فصلی را نیز به این چارچوب افزود. در این پروژه
مؤلفه‌ی فصلی پیاده‌سازی نشده است، چون افق پیش‌بینی مورد نیاز (حدود ۴۵ ثانیه) بسیار کوتاه‌تر
از دوره‌ی هر الگوی فصلی معناداری است.

\subsection{روش‌های سنگین‌تر و چرایی به‌کار نگرفتن آن‌ها}
\label{sec:bg-arima}

خانواده‌ی \lr{ARIMA}~\cite{boxjenkins} و شبکه‌های عصبی بازگشتی مانند \lr{LSTM} در ادبیات
مقیاس‌گذاری پیش‌بینانه فراوان به کار رفته‌اند. این روش‌ها عمداً در این پروژه پیاده‌سازی نشدند
و دلیل آن، هم‌راستا با یافته‌ی اصلی پروژه است:

\begin{quote}
اگر سیگنالی که پیش‌بینی می‌شود ذاتاً نادرست انتخاب شده باشد، پیش‌بینی‌کننده‌ی دقیق‌تر تنها
ناپایداری را \emph{سریع‌تر} می‌کند. مسئله‌ای که این پروژه به آن می‌پردازد، در لایه‌ی انتخاب
سیگنال است و نه در لایه‌ی دقت پیش‌بینی.
\end{quote}

افزون بر این، افق مورد نیاز کوتاه است و بر پایه‌ی سه گام تاریخچه تصمیم گرفته می‌شود؛ در چنین
افقی، برتری روش‌های سنگین بر هموارسازی نمایی محدود است. با این حال، پیاده‌سازی
پیش‌بینی‌کننده پشت یک واسط انجام شده تا افزودن این روش‌ها بعداً کد سیاست را تغییر ندهد
(بخش~\ref{sec:design-forecasters}).

\section{کنترل پس‌خوردی برای سامانه‌های محاسباتی}
\label{sec:bg-control}

مقیاس‌گذار خودکار یک کنترل‌گر پس‌خوردی است؛ این جمله توصیف استعاری نیست، بلکه توصیف ساختاری
است. هلراشتاین و همکاران~\cite{hellerstein2004} چارچوبی برای همین نگاه ارائه کرده‌اند که
واژگان آن در این نوشتار به کار می‌رود:

\begin{itemize}
  \item \textbf{سامانه‌ی تحت کنترل} ناوگان نمونه‌های سرویس است.
  \item \textbf{متغیر کنترل} تعداد نمونه‌هاست.
  \item \textbf{متغیر اندازه‌گیری‌شده} سنجه‌ی بار است.
  \item \textbf{نقطه‌ی تنظیم}\LTRfootnote{Setpoint} نرخ هدف به ازای هر نمونه است. (نام
  مقیاس‌گذار پیاده‌سازی‌شده در این پروژه، \lr{setpoint}، از همین اصطلاح گرفته شده است.)
\end{itemize}

سه مفهوم از این چارچوب در فصل~\ref{ch:results} به‌طور مستقیم به کار می‌آید.

\textbf{بهره‌ی حلقه.}\LTRfootnote{Loop Gain} اگر خروجی کنترل‌گر بر ورودی خودش اثر بگذارد،
حلقه بسته است و علامت این اثر تعیین‌کننده است. پس‌خور \emph{منفی} خطا را کاهش می‌دهد و به
پایداری می‌انجامد؛ پس‌خور \emph{مثبت} خطا را تقویت می‌کند و به واگرایی یا نوسان می‌انجامد.
یافته‌ی مرکزی این پروژه دقیقاً از همین‌جاست: پیش‌بینی بار هر نمونه، مسیری با بهره‌ی مثبت
می‌سازد.

\textbf{تأخیر.} هر تأخیری در مسیر پس‌خور، حاشیه‌ی فاز را کاهش می‌دهد و تمایل سامانه به نوسان
را افزایش می‌دهد. در بخش~\ref{sec:bg-measurement-lag} دیدیم که این تأخیر در مسیر ما بزرگ است.

\textbf{میراسازی.}\LTRfootnote{Damping} افزودن میراگر، دامنه‌ی نوسان را کم می‌کند بی‌آنکه
لزوماً علت آن را برطرف کند. این تمایز — میان \emph{اصلاح} یک قانون کنترل و \emph{پوشاندن}
پیامدهای آن — ستون فقرات بخش~\ref{sec:results-stabilizer} است.

\section{مرور کارهای پیشین}
\label{sec:bg-relatedwork}

\subsection{مرورهای جامع و رده‌بندی‌ها}
\label{sec:bg-surveys}

دو مرور، فضای مسئله را رده‌بندی کرده‌اند. لوریدو-بوتران و همکاران~\cite{lorido2014} روش‌ها را
به پنج خانواده تقسیم می‌کنند که پیش‌تر نام برده شد، و مقیاس‌گذار این پروژه در دو خانواده‌ی
«قواعد آستانه‌ای» و «تحلیل سری‌های زمانی» جای می‌گیرد. کیو و همکاران~\cite{qu2018} در مرور
خود بر چالش‌های باقی‌مانده تمرکز می‌کنند و از جمله به مسئله‌ی نوسان و تعامل میان لایه‌های
مقیاس‌گذاری اشاره دارند.

نکته‌ای که در هر دو مرور دیده می‌شود و برای این پروژه اهمیت دارد: رده‌بندی‌ها عمدتاً بر پایه‌ی
\emph{روش تصمیم‌گیری} انجام شده‌اند و نه بر پایه‌ی \emph{ساختار حلقه}. یعنی این پرسش که
سیگنال ورودی پیش‌بینی‌کننده درون‌زا است یا برون‌زا، محور رده‌بندی نیست.

\subsection{رویکردهای واکنشی}
\label{sec:bg-rw-reactive}

گاندی و همکاران~\cite{gandhi2012} در سامانه‌ی \lr{AutoScale} نشان دادند که بخش بزرگی از
دشواری مدیریت ظرفیت، نه در پیش‌بینی بار بلکه در مدیریت تأخیر راه‌اندازی و پرهیز از واکنش
بیش‌ازحد است. راهبرد آن‌ها، نگه داشتن ظرفیت مازاد به‌جای پیش‌بینی دقیق، از نظر مفهومی به
سازوکار پایدارسازی این پروژه نزدیک است.

نگوین و همکاران~\cite{nguyen2020} \lr{HPA} را به‌صورت تجربی بررسی کرده‌اند و اثر انتخاب
سنجه — منبعی در برابر سفارشی از \lr{Prometheus} — را بر رفتار آن سنجیده‌اند. این کار
نزدیک‌ترین مرجع به بازوهای \lr{hpa-cpu} و \lr{hpa-custom} در ارزیابی این پروژه است.

\subsection{رویکردهای پیش‌بینانه}
\label{sec:bg-rw-predictive}

روی و همکاران~\cite{roy2011} یکی از نخستین کارهای شاخص در استفاده از پیش‌بینی بار برای
تخصیص پیشاپیش منابع‌اند. توکا و همکاران~\cite{toka2021} چند پیش‌بینی‌کننده‌ی مبتنی بر
یادگیری ماشین را در یک موتور مقیاس‌گذاری برای \lr{Kubernetes} به کار برده‌اند و — نکته‌ای که
با یافته‌های این پروژه هم‌راستاست — مسئله را صریحاً به‌شکل یک مبادله میان بیش‌تأمین منابع و
نقض توافق سطح خدمت صورت‌بندی کرده‌اند، نه به‌شکل برتری مطلق.

در صنعت، \lr{Netflix} سامانه‌ی \lr{Scryer}~\cite{scryer} را برای مقیاس‌گذاری پیش‌بینانه
معرفی کرد که بر پایه‌ی الگوهای تکرارشونده‌ی ترافیک، ظرفیت را پیش از نیاز فراهم می‌کند. این
گزارش صنعتی به‌روشنی نشان می‌دهد که سود پیش‌بینی در بارهای \emph{دوره‌ای} متمرکز است — نکته‌ای
که ارزیابی این پروژه به‌صورت مستقل به آن می‌رسد.

\subsection{سامانه‌های صنعتی در مقیاس بزرگ}
\label{sec:bg-rw-industrial}

سامانه‌ی \lr{Borg}~\cite{verma2015} و پس از آن \lr{Autopilot}~\cite{rzadca2020} در گوگل، مسئله
را در مقیاسی متفاوت حل می‌کنند. نکته‌ی مهم \lr{Autopilot} برای این نوشتار، تأکید آن بر
پایداری تصمیم‌ها و پرهیز از تغییرات مکرر است: در مقیاس بزرگ، یک کنترل‌گر پرنوسان حتی اگر
به‌طور میانگین درست تصمیم بگیرد، به دلیل هزینه‌ی خود جابه‌جایی‌ها زیان‌آور است. این دقیقاً
همان نگرانی‌ای است که نیازمندی غیرکارکردی «پرهیز از نوسان» در پیشنهاد این پروژه بیان می‌کند.

\subsection{روش‌شناسی ارزیابی مقیاس‌گذارها}
\label{sec:bg-rw-evaluation}

دسته‌ی چهارم کارهای مرتبط، به \emph{چگونگی سنجش} یک مقیاس‌گذار می‌پردازد. ایلیوشکین و
همکاران~\cite{ilyushkin2018} هفت مقیاس‌گذار را با مجموعه‌ای از سنجه‌های کشسانی مقایسه کرده‌اند
و نشان داده‌اند که رتبه‌بندی مقیاس‌گذارها به انتخاب سنجه به‌شدت حساس است. هربست و
همکاران~\cite{herbst2013} نیز چارچوب مفهومی سنجش کشسانی را ارائه کرده‌اند.

این دسته از کارها پشتوانه‌ی یکی از تصمیم‌های روش‌شناختی این پروژه است: در
فصل~\ref{ch:results} هرگاه یک ستون از جدول نتایج میان بازوها تفکیک نکند، این موضوع صریحاً
گزارش می‌شود و ستون دیگری به‌جای آن انتخاب نمی‌شود بدون آنکه انتخاب اعلام شود.

\subsection{جمع‌بندی مقایسه‌ای}
\label{sec:bg-rw-summary}

جدول~\ref{tab:bg-relatedwork} کارهای مرور شده را از منظر سه پرسش این پروژه کنار هم می‌گذارد.

\begin{table}[H]
\centering
\small
\caption{جایگاه کارهای پیشین نسبت به پرسش‌های این پروژه}
\label{tab:bg-relatedwork}
\begin{tabular}{|l|c|c|c|}
\hline
\textbf{کار} & \textbf{سیاست} & \textbf{بستر} & \textbf{تحلیل پایداری حلقه} \\
\hline
\lr{HPA}~\cite{k8sdocs-hpa} & واکنشی & \lr{Kubernetes} & جزئی (میراگر) \\
\lr{KEDA}~\cite{kedagraduation} & واکنشی & \lr{Kubernetes} & خیر \\
گاندی و همکاران~\cite{gandhi2012} & واکنشی & مرکز داده & جزئی \\
نگوین و همکاران~\cite{nguyen2020} & واکنشی & \lr{Kubernetes} & خیر \\
روی و همکاران~\cite{roy2011} & پیش‌بینانه & ابر & خیر \\
توکا و همکاران~\cite{toka2021} & پیش‌بینانه & \lr{Kubernetes} & خیر \\
\lr{Scryer}~\cite{scryer} & پیش‌بینانه & ابر & خیر \\
\lr{Autopilot}~\cite{rzadca2020} & ترکیبی & داخلی گوگل & بله (پایداری تصمیم) \\
\hline
\textbf{این پروژه} & \textbf{هر دو} & \lr{Kubernetes} & \textbf{بله (محور کار)} \\
\hline
\end{tabular}
\end{table}

\section{شکاف پژوهشی}
\label{sec:bg-gap}

از مرور بالا سه مشاهده به دست می‌آید که با هم شکاف مورد نظر این پروژه را می‌سازند.

\textbf{مشاهده‌ی نخست: پیش‌بینی‌کننده‌ها عمدتاً خارج از حلقه ارزیابی می‌شوند.} در بیشتر
کارهای پیش‌بینانه، کیفیت پیش‌بینی با معیارهایی مانند خطای مطلق میانگین روی یک سری زمانی
\emph{ثبت‌شده} سنجیده می‌شود. اما در سامانه‌ی مستقر، سری زمانی‌ای که پیش‌بینی‌کننده می‌بیند
تحت تأثیر تصمیم‌های خود کنترل‌گر است. ارزیابی خارج از حلقه، به تعریف نمی‌تواند این اثر را
ببیند.

\textbf{مشاهده‌ی دوم: انتخاب سیگنالِ پیش‌بینی معمولاً موضوع بحث نیست.} در ادبیات، بیشتر بر سر
این بحث می‌شود که \emph{چگونه} پیش‌بینی کنیم و کمتر بر سر اینکه \emph{چه چیزی} را پیش‌بینی
کنیم. تمایز میان یک سیگنال برون‌زا (نرخ ورود درخواست‌ها) و یک سیگنال درون‌زا (بار هر نمونه)
تمایزی است که از منظر دقت پیش‌بینی بی‌اهمیت به نظر می‌رسد اما از منظر پایداری حلقه
تعیین‌کننده است.

\textbf{مشاهده‌ی سوم: سازوکارهای میراسازی، ارزیابی را مخدوش می‌کنند.} پنجره‌ی پایدارسازی در
هر مقیاس‌گذار عملی وجود دارد و همیشه در پیکربندی‌های ارزیابی روشن است. اما اگر میراگر روشن
باشد، سنجه‌های متعارف نمی‌توانند یک قانون کنترل سالم را از یک قانون کنترل معیوبِ میراشده
تفکیک کنند. تا جایی که در این مرور دیده شد، ارزیابی‌های موجود این متغیر را به‌صورت یک بازوی
حذفی مستقل بررسی نمی‌کنند.

بر پایه‌ی این سه مشاهده، سهم این پروژه چنین است: \emph{ارزیابی مقیاس‌گذاری پیش‌بینانه به‌شکل
یک حلقه‌ی کنترل بسته، با سیگنال پیش‌بینی به‌عنوان متغیر مستقل، و با سازوکار میراسازی به‌عنوان
یک بازوی حذفی صریح.}

\section{جمع‌بندی فصل}
\label{sec:bg-summary}

این فصل ابزارها و مفاهیم لازم را معرفی کرد. سه نکته که در فصل‌های بعد بیشترین کاربرد را
دارند، یادآوری می‌شود:

\begin{enumerate}
  \item \lr{HPA} با فرمول~\ref{eq:hpa} کار می‌کند، ناحیه‌ی مرده و پنجره‌ی پایدارسازی دارد،
  و در برابر سنجه‌های غایب محافظه‌کارانه رفتار می‌کند. هر سه ویژگی در ارزیابی نقش دارند.
  \item مسیر اندازه‌گیری تأخیری هم‌مرتبه با زمان راه‌اندازی نمونه دارد. بخشی از سود پیش‌بینی،
  جبران همین تأخیر است.
  \item مقیاس‌گذار یک کنترل‌گر پس‌خوردی است، و پایداری آن به علامت بهره‌ی حلقه بستگی دارد و نه
  به دقت مؤلفه‌های آن.
\end{enumerate}

شکاف پژوهشی شناسایی‌شده — نبود ارزیابی درون‌حلقه‌ای برای انتخاب سیگنال پیش‌بینی، و پوشیده
ماندن ناپایداری زیر میراگر — همان چیزی است که طراحی فصل~\ref{ch:design} و روش‌شناسی
فصل~\ref{ch:method} برای پرداختن به آن ساخته شده‌اند.
